\section{Chiến lược huấn luyện}
\label{sec:training}

\subsection{Two-Phase Fine-tuning Strategy}

PhaBERT-CNN sử dụng chiến lược tinh chỉnh hai giai đoạn để đảm bảo sự hội tụ ổn định và thích ứng hiệu quả của mô hình tiền huấn luyện.

\subsubsection{Phase 1: Warm-up (1 epoch)}

\textbf{Mục tiêu:} Khởi tạo task-specific layers tương thích với DNABERT-2 representations

\textbf{Cấu hình:}
\begin{itemize}
    \item \textbf{Freeze}: DNABERT-2 backbone (117M parameters)
    \item \textbf{Train}: CNN branches, attention pooling, classification head (2.1M parameters)
    \item \textbf{Optimizer}: AdamW
    \item \textbf{Learning rate}: $2 \times 10^{-3}$
    \item \textbf{Weight decay}: $1 \times 10^{-4}$
    \item \textbf{Momentum}: $\beta_1 = 0.9$, $\beta_2 = 0.999$
    \item \textbf{LR schedule}: OneCycleLR với 30\% warm-up, max\_lr = $2 \times 10^{-3}$, div\_factor = 5, final\_div\_factor = 10
\end{itemize}

\textbf{Rationale:}
\begin{itemize}
    \item Task-specific layers được khởi tạo ngẫu nhiên, cần learning rate cao để học nhanh
    \item Freeze DNABERT-2 tránh catastrophic forgetting trong giai đoạn đầu
    \item Sau warm-up, task-specific layers đã học được cách xử lý DNABERT-2 embeddings
\end{itemize}

\subsubsection{Phase 2: Full Fine-tuning (tối đa 10 epochs)}

\textbf{Mục tiêu:} Tinh chỉnh toàn bộ mô hình với discriminative learning rates

\textbf{Cấu hình:}

\textit{DNABERT-2 backbone:}
\begin{itemize}
    \item Learning rate: $1 \times 10^{-5}$ (conservative)
    \item Weight decay: $1 \times 10^{-5}$
\end{itemize}

\textit{Task-specific layers:}
\begin{itemize}
    \item Learning rate: $1 \times 10^{-4}$ (10× cao hơn)
    \item Weight decay: $1 \times 10^{-4}$
\end{itemize}

\textit{Common settings:}
\begin{itemize}
    \item Optimizer: AdamW với $\beta_1 = 0.9$, $\beta_2 = 0.999$, $\epsilon = 10^{-6}$
    \item LR schedule: OneCycleLR với 10\% warm-up, div\_factor = 5, final\_div\_factor = 50
    \item Early stopping: Patience = 3 epochs trên validation accuracy
\end{itemize}

\textbf{Rationale:}
\begin{itemize}
    \item DNABERT-2 đã học general genomic patterns, chỉ cần điều chỉnh nhẹ
    \item Task-specific layers cần learning rate cao hơn để học phage lifestyle-specific features
    \item Discriminative learning rates cân bằng giữa adaptation và preservation
\end{itemize}

\subsection{Loss Function}

Binary Cross-Entropy Loss:

\begin{equation}
\mathcal{L}_{\text{BCE}} = -\frac{1}{N}\sum_{i=1}^{N} \left[y_i \log(\hat{y}_i) + (1-y_i)\log(1-\hat{y}_i)\right]
\end{equation}

với:
\begin{itemize}
    \item $y_i \in \{0, 1\}$: True label (0 = temperate, 1 = virulent)
    \item $\hat{y}_i = \sigma(\text{logits}_i) \in (0, 1)$: Predicted probability
    \item $N$: Batch size
\end{itemize}

\subsection{Training Configuration}

\subsubsection{Batch size và gradient accumulation}

\begin{itemize}
    \item \textbf{Effective batch size}: 32
    \item \textbf{Physical batch size}: 8 (giới hạn GPU memory)
    \item \textbf{Gradient accumulation steps}: 4
\end{itemize}

\subsubsection{Mixed precision training}

Sử dụng Automatic Mixed Precision (AMP) để:
\begin{itemize}
    \item Giảm memory consumption
    \item Tăng tốc training ~2×
    \item Duy trì numerical stability với loss scaling
\end{itemize}

\subsubsection{Hardware}

\begin{itemize}
    \item \textbf{GPU}: NVIDIA RTX A4000 16GB
    \item \textbf{CPU}: 32 cores
    \item \textbf{RAM}: 64GB
    \item \textbf{Training time}: ~2-3 giờ/fold (warm-up + full fine-tuning)
\end{itemize}

\subsection{Regularization}

\begin{itemize}
    \item \textbf{Dropout}: 0.1 trong attention pooling và classification head
    \item \textbf{Weight decay}: AdamW với decoupled weight decay
    \item \textbf{Early stopping}: Dừng khi validation accuracy không cải thiện sau 3 epochs
    \item \textbf{Gradient clipping}: Max norm = 1.0 để tránh exploding gradients
\end{itemize}

\subsection{Model Selection}

\begin{itemize}
    \item \textbf{Metric}: Validation accuracy
    \item \textbf{Strategy}: Lưu checkpoint có validation accuracy cao nhất
    \item \textbf{Final model}: Trung bình ensemble của 5 models từ 5 folds (optional)
\end{itemize}

\subsection{Reproducibility}

Để đảm bảo reproducibility:
\begin{itemize}
    \item Random seed: 42 cho tất cả random number generators
    \item Deterministic algorithms: torch.use\_deterministic\_algorithms(True)
    \item CUDA deterministic: torch.backends.cudnn.deterministic = True
    \item Stratified splits: Sử dụng cùng random seed cho cross-validation
\end{itemize}

\subsection{Tổng kết quy trình huấn luyện}

\begin{algorithm}[h]
\caption{PhaBERT-CNN Training Pipeline}
\label{alg:training}
\begin{algorithmic}[1]
\State \textbf{Input:} Dataset $\mathcal{D}$, pre-trained DNABERT-2
\State \textbf{Output:} Fine-tuned PhaBERT-CNN model
\For{each fold $k$ in 5-fold CV}
    \State Split $\mathcal{D}$ into $\mathcal{D}_{\text{train}}^{(k)}$, $\mathcal{D}_{\text{val}}^{(k)}$, $\mathcal{D}_{\text{test}}^{(k)}$
    \State Apply undersampling on $\mathcal{D}_{\text{train}}^{(k)}$
    \State Initialize PhaBERT-CNN with pre-trained DNABERT-2
    \State \textbf{// Phase 1: Warm-up}
    \State Freeze DNABERT-2 backbone
    \State Train task-specific layers for 1 epoch with LR = $2 \times 10^{-3}$
    \State \textbf{// Phase 2: Full Fine-tuning}
    \State Unfreeze all parameters
    \State Train with discriminative LRs (DNABERT-2: $1 \times 10^{-5}$, others: $1 \times 10^{-4}$)
    \State Apply early stopping with patience = 3
    \State Save best model based on validation accuracy
    \State Evaluate on $\mathcal{D}_{\text{test}}^{(k)}$
\EndFor
\State Aggregate results across 5 folds
\State \Return Best models from each fold
\end{algorithmic}
\end{algorithm}

Chiến lược huấn luyện này đảm bảo PhaBERT-CNN học hiệu quả từ DNABERT-2 pre-training trong khi thích ứng tốt với nhiệm vụ phân loại lối sống phage.
